
%----------------------------------------------------------------------------------------
%	DOCUMENT DEFINITION
%----------------------------------------------------------------------------------------

% article class because we want to fully customize the page and not use a cv template
\documentclass[a4paper,12pt]{article}

%----------------------------------------------------------------------------------------
%	FONT
%----------------------------------------------------------------------------------------

% % fontspec allows you to use TTF/OTF fonts directly
% \usepackage{fontspec}
% \defaultfontfeatures{Ligatures=TeX}

% % modified for ShareLaTeX use
% \setmainfont[
% SmallCapsFont = Fontin-SmallCaps.otf,
% BoldFont = Fontin-Bold.otf,
% ItalicFont = Fontin-Italic.otf
% ]
% {Fontin.otf}

%----------------------------------------------------------------------------------------
%	PACKAGES
%----------------------------------------------------------------------------------------
\usepackage{url}
\usepackage{parskip} 	

%other packages for formatting
\RequirePackage{color}
\RequirePackage{graphicx}
\usepackage[usenames,dvipsnames]{xcolor}
\usepackage[scale=0.9]{geometry}

%tabularx environment
\usepackage{tabularx}

%for lists within experience section
\usepackage{enumitem}

% centered version of 'X' col. type
\newcolumntype{C}{>{\centering\arraybackslash}X} 

%to prevent spillover of tabular into next pages
\usepackage{supertabular}
\usepackage{tabularx}
\newlength{\fullcollw}
\setlength{\fullcollw}{0.47\textwidth}

%custom \section
\usepackage{titlesec}				
\usepackage{multicol}
\usepackage{multirow}

%CV Sections inspired by: 
%http://stefano.italians.nl/archives/26
\titleformat{\section}{\Large\scshape\raggedright}{}{0em}{}[\titlerule]
\titlespacing{\section}{0pt}{10pt}{10pt}

%for publications
\usepackage[style=authoryear,sorting=ynt, maxbibnames=2]{biblatex}

%Setup hyperref package, and colours for links
\usepackage[unicode, draft=false]{hyperref}
\definecolor{linkcolour}{rgb}{0,0.2,0.6}
\hypersetup{colorlinks,breaklinks,urlcolor=linkcolour,linkcolor=linkcolour}
\addbibresource{citations.bib}
\setlength\bibitemsep{1em}

%for social icons
\usepackage{fontawesome5}

%debug page outer frames
%\usepackage{showframe}


% job listing environments
\newenvironment{jobshort}[2]
    {
    \begin{tabularx}{\linewidth}{@{}l X r@{}}
    \textbf{#1} & \hfill &  #2 \\[3.75pt]
    \end{tabularx}
    }
    {
    }

\newenvironment{joblong}[2]
    {
    \begin{tabularx}{\linewidth}{@{}l X r@{}}
    \textbf{#1} & \hfill &  #2 \\[3.75pt]
    \end{tabularx}
    \begin{minipage}[t]{\linewidth}
    \begin{itemize}[nosep,after=\strut, leftmargin=1em, itemsep=3pt,label=--]
    }
    {
    \end{itemize}
    \end{minipage}    
    }



%----------------------------------------------------------------------------------------
%	BEGIN DOCUMENT
%----------------------------------------------------------------------------------------
\begin{document}

% non-numbered pages
\pagestyle{empty} 

%----------------------------------------------------------------------------------------
%	TITLE
%----------------------------------------------------------------------------------------

% \begin{tabularx}{\linewidth}{ @{}X X@{} }
% \huge{Your Name}\vspace{2pt} & \hfill \emoji{incoming-envelope} email@email.com \\
% \raisebox{-0.05\height}\faGithub\ username \ | \
% \raisebox{-0.00\height}\faLinkedin\ username \ | \ \raisebox{-0.05\height}\faGlobe \ mysite.com  & \hfill \emoji{calling} number
% \end{tabularx}

\begin{tabularx}{\linewidth}{@{} C @{}}
\Huge{Juan Jara} \\[7.5pt]
\href{https://github.com/Juanchumu}{\raisebox{-0.05\height}\faGithub\ Juanchumu} \ $|$ \ 
\href{https://linkedin.com/in/th3ju4nc3r}{\raisebox{-0.05\height}\faLinkedin\ Juan Jara} \ $|$ \ 
%\href{https://mysite.com}{\raisebox{-0.05\height}\faGlobe \ mysite.com} \ $|$ \ 
\href{mailto:ricard9045@gmail.com}{\raisebox{-0.05\height}\faEnvelope \ ricard9045} \ $|$ \ 
%\href{tel:+54}{\raisebox{-0.05\height}\faMobile \ +54 9 11 } \\
\end{tabularx}

%----------------------------------------------------------------------------------------
% EXPERIENCE SECTIONS
%----------------------------------------------------------------------------------------

%Interests/ Keywords/ Summary
\section{Resumen}
Estudiante de Ciencias Biologicas y Estudiante de Ciencias de Datos, me adapto rapido, hablo ingles fluido, me desempeño bien en GNU/Linux \faLinux, resuelvo y enseño a resolver
con herramientas y servicios de código abierto
Promuevo su uso como opción profesional viable y ética, sin dependencia de licencias.

%Experience
\section{Experciencia Laboral}

%\begin{jobshort}{Datos Geospaciales}{Sep 2024 -  Ag 2025}
%long long line of blah blah that will wrap when the table fills the column width long long line of blah blah that will wrap when the table fills the column width long long line of blah blah that will wrap when the table fills the column width long long line of blah blah that will wrap when the table fills the column width
%\end{jobshort}


\begin{joblong}{Datos Geospaciales}{Sep 2024 -  Ag 2025}
\item Carga, edición y validación de datos en
OpenStreetMap.
\item Uso de datos geoespaciales para la
elaboraciónde mapas temáticos mediante
software SIG (QGIS-JOSM).
\item Contribuciones enfocadas en infraestructura
publica, base de datos de productores
agroecologicos del partido de lujan,
divulgación del uso y utilidades open-source.
\end{joblong}

\begin{joblong}{Fauna en Suelo}{Sep 2022 -  Ag 2023}
\item Separacion y Clasificacion de Coleopteros
hasta especie, pitfall.
\item Generacion de una base de datos robusta.
\end{joblong}
  
%Projects
\section{Proyectos}

\begin{tabularx}{\linewidth}{ @{}l r@{} }
\textbf{Aprender2016-2024} & \hfill \href{https://github.com/Juanchumu/Aprender2016-2024}{Link to GitHubRepo} \\[3.75pt]
\multicolumn{2}{@{}X@{}}{Procesamiento de los datos de las encuestas Aprender hechas a estudiantes de secundaria, modelado y generacion de un Front-End, para producir una herramienta para predecir el desempeño de los grupos.}  \\

\textbf{Predicion de Valor de Casas} & \hfill \href{https://github.com/Juanchumu/TpLabDatos}{Link to GitHubRepo} \\[3.75pt]

\multicolumn{2}{@{}X@{}}{Utilizando una base de datos de casas, se entrenaron modelos de ML, se genero un frontend utilizando streamlit para que un usuario haga predicciones.}  \\


\textbf{Saboteadores} & \hfill \href{https://github.com/Juanchumu/saboteadores}{Link to GitHubRepo} \\[3.75pt]

\multicolumn{2}{@{}X@{}}{Juego Multijugador basado en el Juego de Cartas Saboteur, hecho en java, utilizando javaFX, con interfaz de consola e interfaz Grafica.}  \\


\textbf{EscuelasMapRoulette} & \hfill \href{https://github.com/Juanchumu/EscuelasMapRoulette}{Link to GitHubRepo} \\[3.75pt]

\multicolumn{2}{@{}X@{}}{
Este proyecto tiene como objetivo procesar y transformar el conjunto de datos geográficos del Mapa Escolar, con el fin de integrarlos en el ecosistema de OpenStreetMap (OSM) mediante la plataforma MapRoulette.

El flujo de trabajo incluye la limpieza y estandarización del dataset original, el enriquecimiento de los registros con etiquetas (tags) compatibles con el esquema de OSM (como amenity=school, isced:level, operator:type, entre otros), y la generación de una capa geoespacial apta para su análisis y edición colaborativa.

Posteriormente, los datos son convertidos a un formato GeoJSON que puede ser cargado en MapRoulette, permitiendo a la comunidad validar, corregir y completar los datos directamente sobre el mapa, respetando las convenciones del modelado en OSM.

De esta manera, se busca mejorar la cobertura, precisión y semántica de la información educativa en OpenStreetMap, especialmente en zonas donde los datos están incompletos o desactualizados.
	}  \\
\end{tabularx}

%----------------------------------------------------------------------------------------
%	EDUCATION
%----------------------------------------------------------------------------------------
\section{Educación}
\begin{tabularx}{\linewidth}{@{}l X@{}}	
2024 - present & Analista Universitario en Ciencias de Datos en \textbf{Universidad Nacional de Luján}  \\
2015 - present & Licenciatura en Ciencias Biologicas en \textbf{Universidad Nacional de Luján} \\

2021 & IEEE - ITBA Python \\

2020 & Titulo Marinero de Puente en \textbf{efocapemm} \hfill \\

2019 & Posgrado de Deteccion de Radiacion nuclear en \textbf{Universidad Nacional de Luján} \\
\end{tabularx}

%----------------------------------------------------------------------------------------
%	PUBLICATIONS
%----------------------------------------------------------------------------------------


%----------------------------------------------------------------------------------------
%	SKILLS
%----------------------------------------------------------------------------------------
\section{Habilidades}
\begin{tabularx}{\linewidth}{@{}l X@{}}

\textbf{Frontend} \\
Angular & \normalsize{Desarrollo de aplicaciones frontend con framework Angular}\\
JavaScript & \normalsize{Programación en JS moderno (ES6+), lógica y manipulación del DOM}\\
Streamlit & \normalsize{Desarrollo de aplicaciones web interactivas para visualización de datos}\\
TypeScript & \normalsize{Desarrollo en JavaScript tipado para aplicaciones escalables}\\
Power BI & \normalsize{Visualización de datos y creación de dashboards interactivos}\\

\textbf{Backend} \\
Node.js & \normalsize{Desarrollo backend y APIs con entorno Node.js}\\
C & \normalsize{Programación de bajo nivel, manejo de memoria y desarrollo de sistemas}\\
Java & \normalsize{Desarrollo de aplicaciones orientadas a objetos y backend}\\
Linux & \normalsize{Administración de sistemas y uso de entornos Linux}\\
Ubuntu Server & \normalsize{Configuración y administración de servidores basados en Ubuntu}\\
SSH & \normalsize{Acceso remoto seguro y administración de servidores}\\
Docker & \normalsize{Contenerización de aplicaciones y gestión de entornos}\\
FastAPI & \normalsize{Desarrollo de APIs rápidas y modernas en Python}\\

\textbf{Data Science} \\
Python & \normalsize{Programación general, análisis de datos y scripting}\\
R & \normalsize{Análisis estadístico y procesamiento de datos}\\
R Studio & \normalsize{Uso del entorno RStudio para análisis de datos}\\
JupyterLab & \normalsize{Notebooks interactivos para análisis y prototipado}\\
Google Colab & \normalsize{Ejecución de notebooks en la nube con soporte GPU}\\

\textbf{Machine Learning} \\
LightGBM & \normalsize{Modelos de machine learning basados en gradient boosting eficientes}\\
XGBoost & \normalsize{Implementación optimizada de gradient boosting para modelos predictivos}\\
Redes neuronales (Perceptrón) & \normalsize{Modelos básicos de redes neuronales para clasificación y regresión}\\
Random Forest & \normalsize{Algoritmos de ensamble basados en árboles de decisión}\\

\textbf{Herramientas} \\
Git & \normalsize{Control de versiones distribuido}\\
GitHub & \normalsize{Gestión de repositorios y colaboración en proyectos}\\
LibreOffice & \normalsize{Suite ofimática para documentos, hojas de cálculo y presentaciones}\\
Vim & \normalsize{Edición de texto avanzada en entornos de desarrollo}\\
Arduino & \normalsize{Uso de Arduino como datalogger para adquisición y registro de datos de sensores}\\

Jira & \normalsize{Gestión ágil de proyectos, seguimiento de issues y metodologías Scrum/Kanban}\\
GitHub Projects & \normalsize{Planificación y seguimiento de tareas integradas a repositorios GitHub}\\


\textbf{Bases de datos} \\
SQL & \normalsize{Consultas, modelado y gestión de bases de datos}\\
NoSQL & \normalsize{Diseño y uso de bases de datos no relacionales}\\
Firebird SQL & \normalsize{Administración y consultas en bases de datos Firebird}\\
MySQL & \normalsize{Gestión de bases de datos relacionales y consultas SQL}\\
MongoDB & \normalsize{Base de datos NoSQL orientada a documentos}\\




\end{tabularx}




\vfill
\center{\footnotesize Last updated: \today}

\end{document}
